% Options for packages loaded elsewhere
\PassOptionsToPackage{unicode}{hyperref}
\PassOptionsToPackage{hyphens}{url}
%
\documentclass[
  11pt,
]{article}
\usepackage{amsmath,amssymb}
\usepackage{iftex}
\ifPDFTeX
  \usepackage[T1]{fontenc}
  \usepackage[utf8]{inputenc}
  \usepackage{textcomp} % provide euro and other symbols
\else % if luatex or xetex
  \usepackage{unicode-math} % this also loads fontspec
  \defaultfontfeatures{Scale=MatchLowercase}
  \defaultfontfeatures[\rmfamily]{Ligatures=TeX,Scale=1}
\fi
\usepackage{lmodern}
\ifPDFTeX\else
  % xetex/luatex font selection
\fi
% Use upquote if available, for straight quotes in verbatim environments
\IfFileExists{upquote.sty}{\usepackage{upquote}}{}
\IfFileExists{microtype.sty}{% use microtype if available
  \usepackage[]{microtype}
  \UseMicrotypeSet[protrusion]{basicmath} % disable protrusion for tt fonts
}{}
\makeatletter
\@ifundefined{KOMAClassName}{% if non-KOMA class
  \IfFileExists{parskip.sty}{%
    \usepackage{parskip}
  }{% else
    \setlength{\parindent}{0pt}
    \setlength{\parskip}{6pt plus 2pt minus 1pt}}
}{% if KOMA class
  \KOMAoptions{parskip=half}}
\makeatother
\usepackage{xcolor}
\usepackage[margin=25mm]{geometry}
\usepackage{color}
\usepackage{fancyvrb}
\newcommand{\VerbBar}{|}
\newcommand{\VERB}{\Verb[commandchars=\\\{\}]}
\DefineVerbatimEnvironment{Highlighting}{Verbatim}{commandchars=\\\{\}}
% Add ',fontsize=\small' for more characters per line
\newenvironment{Shaded}{}{}
\newcommand{\AlertTok}[1]{\textcolor[rgb]{1.00,0.00,0.00}{\textbf{#1}}}
\newcommand{\AnnotationTok}[1]{\textcolor[rgb]{0.38,0.63,0.69}{\textbf{\textit{#1}}}}
\newcommand{\AttributeTok}[1]{\textcolor[rgb]{0.49,0.56,0.16}{#1}}
\newcommand{\BaseNTok}[1]{\textcolor[rgb]{0.25,0.63,0.44}{#1}}
\newcommand{\BuiltInTok}[1]{\textcolor[rgb]{0.00,0.50,0.00}{#1}}
\newcommand{\CharTok}[1]{\textcolor[rgb]{0.25,0.44,0.63}{#1}}
\newcommand{\CommentTok}[1]{\textcolor[rgb]{0.38,0.63,0.69}{\textit{#1}}}
\newcommand{\CommentVarTok}[1]{\textcolor[rgb]{0.38,0.63,0.69}{\textbf{\textit{#1}}}}
\newcommand{\ConstantTok}[1]{\textcolor[rgb]{0.53,0.00,0.00}{#1}}
\newcommand{\ControlFlowTok}[1]{\textcolor[rgb]{0.00,0.44,0.13}{\textbf{#1}}}
\newcommand{\DataTypeTok}[1]{\textcolor[rgb]{0.56,0.13,0.00}{#1}}
\newcommand{\DecValTok}[1]{\textcolor[rgb]{0.25,0.63,0.44}{#1}}
\newcommand{\DocumentationTok}[1]{\textcolor[rgb]{0.73,0.13,0.13}{\textit{#1}}}
\newcommand{\ErrorTok}[1]{\textcolor[rgb]{1.00,0.00,0.00}{\textbf{#1}}}
\newcommand{\ExtensionTok}[1]{#1}
\newcommand{\FloatTok}[1]{\textcolor[rgb]{0.25,0.63,0.44}{#1}}
\newcommand{\FunctionTok}[1]{\textcolor[rgb]{0.02,0.16,0.49}{#1}}
\newcommand{\ImportTok}[1]{\textcolor[rgb]{0.00,0.50,0.00}{\textbf{#1}}}
\newcommand{\InformationTok}[1]{\textcolor[rgb]{0.38,0.63,0.69}{\textbf{\textit{#1}}}}
\newcommand{\KeywordTok}[1]{\textcolor[rgb]{0.00,0.44,0.13}{\textbf{#1}}}
\newcommand{\NormalTok}[1]{#1}
\newcommand{\OperatorTok}[1]{\textcolor[rgb]{0.40,0.40,0.40}{#1}}
\newcommand{\OtherTok}[1]{\textcolor[rgb]{0.00,0.44,0.13}{#1}}
\newcommand{\PreprocessorTok}[1]{\textcolor[rgb]{0.74,0.48,0.00}{#1}}
\newcommand{\RegionMarkerTok}[1]{#1}
\newcommand{\SpecialCharTok}[1]{\textcolor[rgb]{0.25,0.44,0.63}{#1}}
\newcommand{\SpecialStringTok}[1]{\textcolor[rgb]{0.73,0.40,0.53}{#1}}
\newcommand{\StringTok}[1]{\textcolor[rgb]{0.25,0.44,0.63}{#1}}
\newcommand{\VariableTok}[1]{\textcolor[rgb]{0.10,0.09,0.49}{#1}}
\newcommand{\VerbatimStringTok}[1]{\textcolor[rgb]{0.25,0.44,0.63}{#1}}
\newcommand{\WarningTok}[1]{\textcolor[rgb]{0.38,0.63,0.69}{\textbf{\textit{#1}}}}
\setlength{\emergencystretch}{3em} % prevent overfull lines
\providecommand{\tightlist}{%
  \setlength{\itemsep}{0pt}\setlength{\parskip}{0pt}}
\setcounter{secnumdepth}{-\maxdimen} % remove section numbering
\ifLuaTeX
  \usepackage{selnolig}  % disable illegal ligatures
\fi
\usepackage{bookmark}
\IfFileExists{xurl.sty}{\usepackage{xurl}}{} % add URL line breaks if available
\urlstyle{same}
\hypersetup{
  hidelinks,
  pdfcreator={LaTeX via pandoc}}

\author{}
\date{}

\begin{document}

\section{Simultaneous model recovery and preservation of a singularity
endpoint}\label{simultaneous-model-recovery-and-preservation-of-a-singularity-endpoint}

Research draft, 6 October 2026. \textbf{Status: PARTIAL;
candidate=false.}

The entire unchanged residual P19.2 existence problem is not solved
here. We reconstruct the two supplied candidates, prove an
endpoint-preserving construction for an explicitly restricted
infinite-chain input, and give an analytic warning about changing a
full-mass initial potential. The restricted construction is not promoted
to a theorem for arbitrary log-discrepancy models, arbitrary prescribed
singularities, or arbitrary compact Kähler manifolds. Historical review
descriptions in the supplied documents are not proofs. The sole
substantial theorem grant used below is the complete supplied P19.1
mathematical theorem in its exact signature.

\subsection{1. Original target and
conventions}\label{original-target-and-conventions}

Original P19.2 asks for model approximations in the residual hybrid
programme preserving the required energy, entropy, and singularity data.
The dossier expressly excludes energy recovery alone and retains the
published projective Fujita regimes as established subcases. Its
completion criterion requires the entire claimed residual signature; a
conditional mechanism or a narrower theorem cannot close the atom.

The dossier does not specify a universal numerical list of prescribed
levels or one universal endpoint topology. The additional assignment
requires finer P data, including endpoints. We keep separate the
following properties: equality of prescribed envelopes at each index;
eventual equality at a fixed parameter; convergence of envelopes;
convergence in the singularity pseudometric; convergence of every mixed
raw mass; and preservation of exact polar sets. None is silently
identified with another.

Fix a connected smooth compact Kähler manifold \((X,\omega)\), of
dimension \(n\geq1\), and write \[
 \alpha=[\omega],\qquad V=\int_X\omega^n,\qquad
 \mu_0=\omega^n/V,\qquad dd^c=\frac{i}{2\pi}\partial\bar\partial,
 \qquad t=-\log|\tau|^2.
\] In particular \(dd^c\log|f|^2=[\operatorname{div}(f)]\). Use
probability Monge--Ampère measure and energy \[
 \operatorname{MA}(u)=V^{-1}\langle\omega_u^n\rangle,
 \quad E(u)=\frac{1}{(n+1)V}\sum_{k=0}^n
              \int_Xu\,\omega_u^k\wedge\omega^{n-k},
 \quad \omega_u=\omega+dd^cu.
\] On \(\mathcal E^1(\omega)\), \(E(0)=0\), \(E(u+C)=E(u)+C\), and these
measures have mass one. Entropy is the full probability relative entropy
\(H(u)=\operatorname{Ent}(\operatorname{MA}(u)\mid\mu_0)\), with value
\(+\infty\) when absolute continuity fails. Non-Archimedean data live on
the original \(X^\beth\) over trivially valued \(\mathbf C\), including
semivaluations with nontrivial support. We never discard those supports.

For a smooth equivariant model \(\rho:\mathcal Y\to X\times\mathbf P^1\)
and a vertical real divisor \(D\), the datum is \(f_D(v_F)=D_F/b_F\),
\(b_F=\operatorname{ord}_F(\tau)\), and compatibility means
\(U+\log|s_D|^2=O(1)\). An actual approximating model has
\(\rho_X^*\alpha+D\) relatively Kähler. In the algebraic examples a
rational \(D\) gives an actual relatively ample rational line bundle.
Write \[
 H_\alpha(\phi)=\int_{X^\beth}A_X\operatorname{MA}_\alpha(\phi),
 \qquad A_X(v_{\rm triv})=0.
\] The measure in this formula is a probability. At a Gauss component
its discrepancy over \(X\) is \((\kappa_F+1-b_F)/b_F\), where
\(\kappa_F\) is the relative canonical coefficient over
\(X\times\mathbf P^1\).

For a ray \(u\) define \(\chi_u(\lambda)=\inf_{t\geq0}(u_t-t\lambda)\).
Its level mass is
\(V^{-1}\int\langle\omega_{\chi_u(\lambda)}^n\rangle\); an identically
\(-\infty\) level is assigned mass zero. These are raw masses, not
probability measures. No zero-mass level is divided by its mass. The
envelope is
\[ P_\omega[u]=\left(\lim_{C\to\infty}P_\omega(u+C,0)\right)^*.\] P
equivalence refers to this finer ceiling/envelope data. Multiplier-ideal
I equivalence is not a substitute.

\subsection{2. What the primary log-discrepancy-model regularization
states}\label{what-the-primary-log-discrepancy-model-regularization-states}

We read the primary Darvas--Zhang text, \emph{A YTD correspondence for
constant scalar curvature metrics}, arXiv:2509.15173v2, formula (11),
Section 9, and the complete proof of Theorem 10.1. In its polarized
signature a smooth dominating test configuration
\((\mathcal X,\mathcal L)\) has, for \(\beta\in\mathbf Q_{>1}\),
log-discrepancy model \[
 \mathcal L_\beta=\mathcal L+\beta^{-1}C_{\mathcal X},\qquad
 C_{\mathcal X}=K_{\mathcal X/(X\times\mathbf P^1)}
                          +\mathcal X_{0,\rm red}-\mathcal X_0.
\] Its NA datum is the psh envelope \[
 P_\alpha\bigl(\phi_{\mathcal L}
                  +\beta^{-1}A_X\circ r_{\mathcal X}\bigr),
\] not in general the PL obstacle itself. The rational bundle need not
be relatively semiample. Theorem 10.1 assumes a polarized manifold and a
finite-energy ray of finite radial entropy. It produces smooth test
configurations and \(\beta_k\to\infty\) whose log-discrepancy-model rays
converge in chordal \(d_1^c\) and radial entropy. Here
\(d_1^c(u,v)=\lim_{t\to\infty}d_1(u_t,v_t)/t\).

The proof first fixes \(\beta\), uses multiplier-ideal regularization
and continuity of radial \(\beta\) entropy, then replaces the test
configuration by its log-discrepancy model. Its order and energy
estimate give chordal error at most
\((\operatorname{Ent}(u)+\beta^{-1}+1)/\beta\); its entropy inequality
and lower semicontinuity give entropy recovery. Lemma 9.1 checks that
the envelope is independent of a further smooth resolution: the new
log-canonical error is effective exceptional and does not change the
translated big volume. Lemma 9.3 keeps all central multiplicities in its
exponential-integrability comparison. No step in these pertinent units
estimates a prescribed endpoint P envelope or asserts exact poles.

This primary theorem and proof are used here to identify the signature,
not as a new grant or as a proof for a transcendental class. The
supplied P19.1 theorem provides actual relatively Kähler recovery in a
fixed smooth compact Kähler class. Neither result, as stated, supplies
the missing endpoint conclusion. The dossier's dated Trusiani polarized
claims are not relisted as wholly open, and neither the static cscK
implication nor any genuinely degenerating-family or automorphism
extension is reproved.

\subsection{3. Exact mathematical inputs used in the
proofs}\label{exact-mathematical-inputs-used-in-the-proofs}

The supplied complete P19.1 theorem, denoted R, is used in precisely the
following fixed smooth compact Kähler signature. For
\(w\in\mathcal E^1(\omega)\) and
\(\phi\in\mathcal E^{1,\rm NA}(\alpha)\) there is an actual maximal weak
ray \(r(w,\phi)\), ordered in both arguments, with
\[ E(r(w,\phi)_t)=E(w)+tE_\alpha(\phi). \tag{3.1}\] Its usual datum is
the actual total-potential datum, not an assigned valuation function. A
maximal ray with a finite-entropy slice has entropy slope
\(H_\alpha(\phi)\). If this is finite, every initial finite-entropy ray
in this branch has finite entropy at finite times. For finite
\(H_\alpha(\phi)\), R supplies actual smooth equivariant relatively
Kähler models \(\eta_j\) with strong convergence to \(\phi\) and
\(H_\alpha(\eta_j)\to H_\alpha(\phi)\). It also supplies calibrated
mixed slopes, fixed log-canonical tests, and joint continuity of the
maximal-ray construction. No additional P-endpoint theorem is contained
in this grant.

We use the following finite-energy foundations in the domains just
stated. Ordered \(d_1(a,b)=E(a)-E(b)\) for \(a\geq b\); the supporting
inequality is \(E(a)-E(b)\leq\int(a-b)\operatorname{MA}(b)\); energy is
continuous in \(d_1\). Strong NA convergence gives energy and supremum
convergence and \(\int|\eta_j-\phi|\operatorname{MA}_\alpha(\psi)\to0\)
for fixed \(\psi\in\mathcal E^{1,\rm NA}\). These are the foundations
explicitly retained in R; the pertinent primary MPWN Lemma 5.0.5 and
mixed estimates were also checked. All measures in these uses have mass
one. The NA psh class is closed under uniform limits; bounded psh data
have finite energy, and uniform bounded convergence gives energy and
strong convergence. These are the bounded part of the same psh/energy
foundations, used when identifying the explicit infinite-chain limit.

Mesquita--Piccione, arXiv:2409.06221v2, Lemma 5.3.2, gives a mixed
energy-pairing difference for mean-zero \(\mathcal E^1(\omega)\) slots
bounded by \(C\varepsilon^q R_0^{1-q}\), \(q=2^{-n}\). Here
\(\varepsilon\) bounds the relative, constant-invariant \(J\) distances
of replaced slots, \(R_0\) bounds their \(J\) values plus the compact
psh supremum-minus-mean constant, and \(C\) depends on dimension and the
fixed smooth forms. The source proof extends the smooth inequality by
decreasing finite-energy continuity. Its Theorem 5.3.4 proof splits a
fixed signed PL test into positive weights on one fixed model, and
bounds relative \(J\) by \(d_1\). We use that distance, not a supposed
bound for \(d_1\) after arbitrary mean subtraction. MP's \(i/\pi\)
convention is transported with background \(2\omega\), the same
potentials at \(t=2s\), and datum \(2\phi\) in \(2\alpha\). The grant's
calibration fixes mixed slopes in squared-log time.

BDL, \emph{Convexity of the extended K-energy}, Lemma 3.1 and Theorem
4.7, provide smooth \(d_1\) approximants with entropy convergence and
the lower semicontinuous geodesically convex functional \[
 M=H+G,\qquad G=T_{-\operatorname{Ric}(\omega)}+\bar S E,
 \quad T_\xi(u)=V^{-1}\sum_{k=0}^{n-1}
       \int_Xu\,\xi\wedge\omega_u^k\wedge\omega^{n-1-k},
\] where \(\operatorname{Ric}=-dd^c\log\omega^n\) and
\(\bar S=nV^{-1}\int\operatorname{Ric}\wedge\omega^{n-1}\). The
contracted energies are \(d_1\) continuous. In their
\(i\partial\bar\partial\) convention use potential \(u/(2\pi)\); their
Ricci and scalar curvature scale by \(2\pi\), cancelling the energy
scaling. Full probability entropy retains coefficient one. R supplies
its calibrated maximal-ray mixed slopes.

For the explicit curves we use only the analytic initial-zero Legendre
correspondence: DZ, arXiv:2208.08324v4, Theorem 3.7(ii), proves that a
parameter-concave, usc, P-maximal test curve with initial limit zero has
an actual weak ray as inverse Legendre transform. Their complete Perron
replacement proof was read. Our curves are bounded in the test-curve
sense and satisfy \(-rt\leq u_t\leq0\), so continuity at zero and finite
energy require no general singular-ray completion argument. DXZ Theorem
2.6 and its relative definitions were read for comparison, not applied
to replace a singular initial potential by a smooth background. The
ordered DDNL Lemma 3.4 formula, after operator and probability
transport, gives on our curve
\[ d_{\mathcal S}([w_b],[w_d])=|b-d|/2. \tag{3.2}\] Its full mixed-mass
formula is \((n+1)^{-1}\sum_{k=0}^n(M_k(v)-M_k(u))\) for ordered
\([u]\leq[v]\), where
\(M_k(u)=V^{-1}\int\omega^{n-k}\wedge\langle\omega_u^k\rangle\).

For bounded changes of singularity type we use DDNL, arXiv:1705.05796v3,
Theorem 2.4: for smooth closed real pseudoeffective forms
\(\theta_1,\ldots,\theta_n\), if each
\(u_i\in\operatorname{PSH}(X,\theta_i)\) is less singular than
\(v_i\in\operatorname{PSH}(X,\theta_i)\), then
\(\int_X\langle\theta_{1,u_1}\wedge\cdots\wedge\theta_{n,u_n}\rangle
\geq\int_X\langle\theta_{1,v_1}\wedge\cdots\wedge\theta_{n,v_n}\rangle\).
Its complete primary truncation proof was read. Applying both orders
when the differences are bounded proves equality of all mixed raw
masses. This theorem has no positive-top-mass assumption; it does not
identify P envelopes at zero mass.

Stacks Lemma 33.44.15, Tag 0B5Y, applies to any proper scheme of
dimension at most one over a field: an invertible sheaf is ample exactly
when it has positive degree on every one-dimensional irreducible
component. Tag 09MS proves invariance under a closed Noetherian
thickening with the same underlying space. Tag 0D2S says that a proper
morphism and an invertible sheaf ample on one fiber have a neighborhood
where the sheaf is relatively ample. The first two complete proof units
are supplied; the last complete primary proof was freshly read. These
are applied to actual scheme fibers, including all nonreduced
multiplicities.

\subsection{4. Reconstructing candidate187: singular starts and a joint
entropy
limit}\label{reconstructing-candidate187-singular-starts-and-a-joint-entropy-limit}

\textbf{Lemma 4.1 (uniform initial-point restoration).} Normalize
\(\sup_Xw_j=0\), suppose \(A_j=-E(w_j)\) is bounded, and let
\(\eta_j\to\phi\) strongly. Put \(a_{j,t}=r(0,\eta_j)_t\) and
\(b_{j,t}=r(w_j,\eta_j)_t\). For \(t\geq1\), \[
 a_{j,t}\geq b_{j,t},\quad d_1(a_{j,t},b_{j,t})=A_j,
 \quad |G(a_{j,t})-G(b_{j,t})|
 \leq C\{A_j+A_j^q(t+A_j+1)^{1-q}\}. \tag{4.1}
\] For each fixed divisor-compatible smooth test \(\Psi_t\), \[
 \left|\int\Psi_t\operatorname{MA}(a_{j,t})
             -\int\Psi_t\operatorname{MA}(b_{j,t})\right|
 \leq C_\Psi A_j^q(t+A_j+1)^{1-q}. \tag{4.2}
\] The constants do not depend on the varying models representing
\(\eta_j\).

\emph{Proof.} Order and (3.1) prove the first two assertions. Set
\(m(u)=\int u\,d\mu_0\), \(u^0=u-m(u)\), and \(J(u)=m(u)-E(u)\). From
zero, \(\sup a_{j,t}=t\sup\eta_j\), so \(J(a_{j,t})\leq Ct\) by the
uniform datum energy and supremum bounds. Since \(m(b)\leq m(a)\) and
\(E(a)-E(b)=A_j\), \(J(b)\leq J(a)+A_j\). The relative \(J(a,b)\) is at
most \(A_j\) by its supporting inequality and is invariant under the two
constant changes. Apply the primary mixed estimate first to the
normalized \emph{energy} pairing. It gives \[
 |E(a^0)-E(b^0)|\leq C A_j^q(t+A_j+1)^{1-q}.
\] The identity \(E(a)-E(b)=m(a)-m(b)+E(a^0)-E(b^0)\) then gives
\[ |m(a)-m(b)|\leq A_j+C A_j^q(t+A_j+1)^{1-q}. \tag{4.3}\] Apply the
estimate to each normalized contracted pairing, restore its means using
(4.3), and use \(E(a)-E(b)=A_j\) for the \(\bar S E\) term. This proves
(4.1). Neither boundedness below of \(w_j\) nor a singular background
change is used.

For (4.2), the integral is the Deligne pairing
\((0,\Psi_t)(\omega,u_t)^n/V\). Split the fixed signed test as in the
primary proof, on one fixed common model. Its positive slots have
\(J=O(t+1)\). Replace the changing slots individually. A constant change
in a ray slot contributes the integral of \(dd^c\Psi_t\) against a
closed product on the compact fiber, which is zero. A test-slot constant
cancels between the two probability measures. Thus no mean-restoration
term remains. The fixed forms and model determine \(C_\Psi\). This
proves the lemma.

\textbf{Theorem 4.2 (candidate187 reconstructed in its exact partial
scope).} Let \(H_\alpha(\phi)<\infty\) and let actual models \(\eta_j\)
be supplied by R. If \(w_j\to w\) in analytic \(d_1\) and
\(\sup_j H(w_j)<\infty\), then for every \(j_\ell,t_\ell\to\infty\), \[
 \frac{H(r(w_{j_\ell},\eta_{j_\ell})_{t_\ell})}{t_\ell}
                     \longrightarrow H_\alpha(\phi). \tag{4.4}
\] In particular this holds for a fixed genuine \(w\in\mathcal E^1\)
with \(H(w)<\infty\). It makes no relative endpoint-envelope assertion.

\emph{Proof.} Supremum-normalize the starts. Entropy and measures are
unchanged; (4.4) is unchanged. Bounded \(d_1\) starts have bounded
energies and suprema, so \(A_j\) is bounded. The \(G(w_j)\) are bounded
by \(d_1\) continuity (and \(G\) is constant invariant); hence
\(M(w_j)\) is uniformly bounded above. It is also bounded below, since
\(H\geq0\).

For each fixed \(j\), the actual maximal ray \(b_j=r(w_j,\eta_j)\) has a
finite-entropy slice at zero. R therefore gives its entropy slope
\(h_j=H_\alpha(\eta_j)<\infty\). Lemma 4.1 divided by \(t\) gives its
contracted slope \(g_j=G_\alpha(\eta_j)\) from the zero-start calibrated
slope. Convexity of \(M\) bounds every finite secant by this finite
recession slope: \[ M(b_{j,t})\leq M(w_j)+t(h_j+g_j). \tag{4.5}\] A
finite-time infinite value would force an infinite recession slope, so
the same argument also checks finiteness of each entropy here.

For zero starts use the genuine common upper maximal ray of datum
\(\lambda_j=\max(\eta_j,\phi)\). The supporting inequality gives \[
 B_j=2E_\alpha(\lambda_j)-E_\alpha(\eta_j)-E_\alpha(\phi)\to0,
 \qquad d_1(a_{j,t},a_t)\leq tB_j,\quad a_t=r(0,\phi)_t.
\] Indeed the maximum energy defect is bounded by
\(\int|\eta_j-\phi|\operatorname{MA}_\alpha(\phi)\to0\). No pointwise
maximum of rays is declared geodesic. The fixed-form estimate, first on
normalized energy and then restoring means, yields the useful bound \[
 \frac{|G(a_{j,t})-G(a_t)|}{t}
 \leq C\{|E_\alpha(\eta_j)-E_\alpha(\phi)|
                  +B_j^q(1+t^{-1})^{1-q}\}. \tag{4.6}
\] Here all changing \(J\) values are \(O(t+1)\) and their relative
distances are at most \(CtB_j\). Applying the normalized energy estimate
restores the mean difference by \(t|\Delta E_\alpha|\) plus the same
Hölder error. This explains all constants in (4.6). A fixed PL-test
integral instead has normalized difference at most \(C_\Psi B_j^q\).

Equations (4.1) and (4.6), and the target mixed slope, give the joint
limit \(G(b_{j,t})/t\to G_\alpha(\phi)\). Strong mixed continuity gives
\(g_j\to G_\alpha(\phi)\); R gives \(h_j\to H_\alpha(\phi)\). Divide
(4.5) by \(t\) to prove the joint entropy upper bound.

For the lower bound first fix one dominating SNC model \(\mathcal Y\).
R's log-canonical test has \(f_\mathcal Y=A_X\circ r_\mathcal Y\) and
\(\Psi_t=\log(\mu_{\mathcal Y,t}/\mu_0)\), with
\(\int\mu_{\mathcal Y,t}\leq C_\mathcal Y(1+t)^n\). In local coordinates
\(\tau=\prod z_i^{b_i}\), the independent logarithmic radial variables
occupy a weighted simplex; its dimension is at most \(n\). Thus this
bound retains all \(b_i\), rather than replacing the fiber by its
reduction. Gibbs' inequality gives \[
 H(b_{j,t})\geq\int\Psi_t\operatorname{MA}(b_{j,t})
                         -\log C_\mathcal Y-n\log(1+t).
\] Replace \(b_j\) by \(a_j\) using (4.2), with error divided by \(t\)
tending uniformly to zero because \(A_j\) is bounded and \(q>0\).
Replace \(a_j\) by \(a\) using its fixed-test \(C_\mathcal YB_j^q\)
bound. The calibrated target slope gives the joint lower bound
\(\int f_\mathcal Y\operatorname{MA}_\alpha(\phi)\). Only after this
joint liminf take the directed supremum of fixed SNC tests. R's
discrepancy-integral identity gives \(H_\alpha(\phi)\). This proves
(4.4), without a uniform bound over varying test models.

BDL's initial regularizers are admissible here: truncate the target
probability density, smooth it with entropy convergence, solve the
smooth complex Monge--Ampère equation, use finite-entropy \(d_1\)
compactness and uniqueness, then restore constants. Lemma 3.1 produces
\(w_j\in\mathcal H_\omega\) with \(d_1(w_j,w)\to0\) and
\(H(w_j)\to H(w)\) whenever \(H(w)<\infty\). This is controlled entropy
at the initial time, not a theorem of entropy continuity at every fixed
later time. In general only lower semicontinuity is available there.
With these starts, R's common upper ray also gives bounded-time \(d_1\)
convergence to \(r(w,\phi)\).

For infinite-entropy starts the lemma still holds but (4.5) need not
have a finite intercept. If starts depend on a chosen \(t_j\), the same
proof works under bounded normalized energies and \(M(w_j)/t_j\to0\);
this is a conditional sufficient condition, not automatic from \(d_1\)
convergence.

\subsection{5. Reconstructing candidate188 on actual
blowups}\label{reconstructing-candidate188-on-actual-blowups}

The following finite-chain lemma supplies the diagnostic and our
constructive result. This proof includes the ideal-invertibility and
Legendre steps; it does not infer an actual model from a formal speed.

\textbf{Lemma 5.1 (finite chain).} Fix \(m\geq1\) and positive rational
numbers \(a_1,\ldots,a_m\) with \(c=\sum a_i<1\). There is an actual
smooth equivariant model of \(\mathbf P^1\) and a relatively ample
rational bundle with the following datum, measure, energy, and ray. On
the branch of zero, write \(v=s\operatorname{ord}_0\), allowing
\(s=+\infty\) at the supported semivaluation. Then \[
 \phi(v)=-\sum_i a_i\min(s,1/i),\quad
 \operatorname{MA}_\alpha(\phi)=(1-c)\delta_{v_{\rm triv}}
                     +\sum_i a_i\delta_{i^{-1}\operatorname{ord}_0}.
 \tag{5.1}
\] It is zero on the other branches. Set \(S_i=\sum_{k\leq i}a_k\),
\(S_0=0\), \(v(x)=1/i\) on \((S_{i-1},S_i)\), and \[
 g(x)=\int_x^c v(s)\,ds\ (0\leq x\leq c),\qquad
 g(x)=0\ (c\leq x\leq1),\qquad r=g(0)=\sum_i a_i/i.
\] The actual zero-start maximal ray is \[
 q_t(\rho)=\max_{0\leq x\leq1}\{x\rho-F(x)-tg(x)\},\quad
 u_t=q_t(\log|z|^2)-\log(1+|z|^2),\quad
 F(x)=x\log x+(1-x)\log(1-x).
 \tag{5.2}
\] The endpoint is \(w_c\), where for \(0<b<1\), \[
 w_b(z)=
 \begin{cases}
 b\log|z|^2-F(b)-\log(1+|z|^2),&|z|^2\leq b/(1-b),\\
 0,&|z|^2\geq b/(1-b).
 \end{cases} \tag{5.3}
\] Put \(w_0=0\). Each \(w_b\) is a genuine P-model potential with pole
coefficient \(b\) at zero and raw mass \(1-b\). The complete level curve
is \[
 \chi(\lambda)=
 \begin{cases}
 0,&\lambda\leq-r,\\
 w_{b(\lambda)},&-r<\lambda<0,\quad g(b(\lambda))=-\lambda,\\
 w_c,&\lambda=0,\\
 -\infty,&\lambda>0.
 \end{cases} \tag{5.4}
\] Finally \[
 E_\alpha(\phi)=-\int_0^1g(x)\,dx
 =-\tfrac12\sum_{i,k}a_ia_k\min(1/i,1/k),\quad
 E(u_t)=tE_\alpha(\phi),\quad H_\alpha(\phi)=r,
 \tag{5.5}
\] \[
 H(u_t)=tr+2\sum_i\int_{S_{i-1}}^{S_i}
        \log(1-x+e^{-t/i}x)\,dx,\qquad -2\leq H(u_t)-tr\leq0.
 \tag{5.6}
\]

\emph{Proof of actual geometry and positivity.} Start with
\(\mathbf P^1_z\times\mathbf P^1_\tau\), blow up \((0,0)\), and
repeatedly blow up the intersection of the newest exceptional component
with the strict central component. Call the exceptional components
\(F_i\) and the strict central component \(F_0\). The local fan vectors
are \((1,i)\); consecutive determinants are one, so these are smooth
charts. All centers are invariant smooth points. The fiber, incidence,
and intersection numbers are \[
 B=F_0+\sum_{i=1}^m iF_i,\quad
 F_0,F_m,F_{m-1},\ldots,F_1,\quad
 F_0^2=-m,\quad F_i^2=-2\ (i<m),\quad F_m^2=-1.
\] Adjacent intersections are one and others zero. These formulas follow
inductively by the blowup intersection rule: a new central multiplicity
is the sum \(i+1\) of the two component multiplicities meeting at its
center. In particular \(v_{F_i}=i^{-1}\operatorname{ord}_0\) and
\(b_i=i\).

Put \[
 d_i=-\sum_k a_k\min(1,i/k),\quad d_0=0,\quad
 D=\sum_{i=1}^m d_iF_i,\quad \mathcal L=p_z^*\mathcal O(1)+D.
 \tag{5.7}
\] On this smooth surface every \(F_i\) is Cartier. Clear denominators
to obtain an actual invertible sheaf. For \(i<m\),
\(D.F_i=-2d_i+d_{i-1}+d_{i+1}=a_i/i\), because the discrete second
difference of \(-a_k\min(1,i/k)\) is \(a_k/k\) at \(i=k\) and zero
otherwise. For \(i=m\), \(D.F_m=-d_m+d_{m-1}=a_m/m\). Also
\(D.F_0=d_m=-c\). Thus
\[ \mathcal L.F_0=1-c>0,\qquad \mathcal L.F_i=a_i/i>0. \tag{5.8}\] Other
fibers have degree one. Stacks 0B5Y proves ampleness on the actual
proper central curve; 09MS alternatively justifies passing through its
nilpotent thickening. Both apply after a denominator and do not require
reduced fibers. Properness and 0D2S give relative ampleness. A
sufficiently positive base twist makes the rational bundle ample on the
projective surface, giving a Kähler representative. The divisor
describing the relative model and its normalized datum remain (5.7): the
base twist is an auxiliary positivity device, not an unnoticed central
translation.

\emph{Proof of compatibility and the datum.} Convex duality gives
\[ \max_{0\leq x\leq1}\{-sx-g(x)\}
                  =-\sum_i a_i\min(s,1/i). \tag{5.9}\] For example
minimizers of \(sx+g(x)\) move through the nodes as \(s\) crosses
\(1/i\), so the derivative of its minimum is the sum of \(a_i\) whose
speeds exceed \(s\); integration from \(s=0\) proves (5.9). At \(s=1/k\)
its value is \(d_k/k\) and at \(s=+\infty\) it is \(-r\). There is no
subtraction of two infinite values at this support.

Omit \(F\) in (5.2), and denote the resulting support weight by \(h\).
Since \(-\log2\leq F\leq0\), \(h\leq q\leq h+\log2\). Its nodes include
\(x=0,S_1,\ldots,S_m=c,1\). For a common denominator \(M\), its
monomials in a local \(\mathcal O(M)\) frame are
\[ z^{MS_i}\tau^{Mg(S_i)}\ (0\leq i\leq m),\qquad z^M. \tag{5.10}\] All
exponents are nonnegative integers. The node \(x=0\) gives a pure power
of \(\tau\), so the ideal is a unit off the central fiber. The node
\(x=1\) gives a unit in the infinity frame and removes a possible
horizontal base locus there.

Here is the invertibility check. In the cone between \((1,k)\) and
\((1,k+1)\), node \(S_k\) minimizes the monomial order on both
generating rays; this follows by the consecutive slopes
\(-1/k,-1/(k+1)\) of \(g\). Every other exponent minus its exponent
therefore pairs nonnegatively with both rays. Those rays are a lattice
basis; the quotient monomials are regular in that smooth affine chart,
and the selected generator has quotient one. The first cone between
\((1,0)\) and \((1,1)\) uses \(x=0\); the last cone between \((1,m)\)
and \((0,1)\) uses \(x=c\). These cover the local blowup fan. Away from
the origin the ideal is already invertible. Its coefficient on \(F_k\)
is \(M\min_x(x+kg(x))=-Md_k\), and on \(F_0\) it is zero. Removing this
canonical factor leaves a locally bounded logarithmic maximum. The
comparison with \(q\) proves exactly \(U+\log|s_D|^2=O(1)\) in actual
bundle frames. Thus (5.9) is the actual model datum on the entire
carrier, including supports.

\emph{Proof of the level calculation and the weak-ray property.} At each
finite \(t\), \(F+tg\) is strictly convex, so \(q_t\) has a unique
maximizing \(x\) on the open moment interval. Convex Lagrange duality
for \(g(x)\leq-\lambda\) is strict for \(-r<\lambda<0\) (use \(x=1\)).
It identifies \(\inf_{t\geq0}(q_t-t\lambda)\) with the maximum of
\(x\rho-F(x)\) over \([b(\lambda),1]\). If the unrestricted maximizing
\(x\) exceeds \(b\), this is the Fubini--Study weight; otherwise its
maximum is at \(b\). This proves the interior formula (5.3). At
\(\lambda\leq-r\) the initial time gives zero and \(u_t\geq-rt\) gives
the reverse bound. At \(\lambda>0\), \(u_t\leq0\) gives \(-\infty\). At
\(\lambda=0\), compactness of maximizing \(x\) forces its cluster points
into \(g^{-1}(0)=[c,1]\) as \(t\to\infty\). Evaluating every fixed
\(x\geq c\) gives the reverse limiting bound, hence the endpoint
\(w_c\).

The ordinary curvature of \(w_b\) has atom \(b\delta_0\), equals
\(\omega\) outside the circle \(|z|^2=b/(1-b)\), and has no circle atom,
because first derivatives match. Removing the pole leaves raw mass
\(1-b\). To check the P property directly, let \(h_0\leq0\) and
\(h_0\leq w_b+C\) be any rooftop competitor. Its local subharmonic
weight \(Q=h_0+\log(1+|z|^2)\) has an atom of at least \(b\) at zero.
Subtracting \(b\log|z|^2\) leaves a subharmonic function bounded above
near zero, so it extends there. On the boundary of the disk its value is
at most \(\log(1+B)-b\log B=-F(b)\), \(B=b/(1-b)\). The maximum
principle gives \(h_0\leq w_b\) inside; outside \(h_0\leq0=w_b\). Since
\(w_b\) itself is a competitor, \(P_\omega[w_b]=w_b\).

The level curve (5.4) is parameter-concave and usc as the infimum of
time-affine functions in its definition. It is P-model at every
nontrivial level, including the endpoint, and zero below \(-r\). The
analytic DZ correspondence applies with background \(\omega\) and
initial zero. Convex biconjugacy gives inverse transform (5.2). Thus it
is an actual bounded weak ray, strongly continuous at zero. This
establishes geodesicity without assigning it to a subgeodesic.

\emph{Proof of measure, energy, and entropy.} Its model component masses
are \(b_i\mathcal L.F_i=a_i\) and \(1-c\) on \(F_0\), proving (5.1). The
total-space canonical coefficient on the ray \((1,i)\) is \(i\);
subtracting the time multiplicity gives \(A_X(v_{F_i})=1/i\). Thus its
discrepancy entropy is \(r\), including the nonreduced correction.
Intersecting \(D\) with (5.8) proves the double sum in (5.5), since
\(\min(1,i/k)/i=\min(1/i,1/k)\). Integrating \(g\) proves the same
formula.

The radial probability is \(dx\), where \(x=\partial_\rho q_t\);
intervals where \(x\) is constant have zero curvature mass. Envelope
differentiation gives \(\partial_tu_t=-g(x)\), so the bounded
weak-geodesic energy derivative gives \(E'(u_t)=-\int_0^1g(x)dx\) and
\(E(u_0)=0\). Equality of actual and datum energy slopes implies
maximality for the usual datum by R's ordered energy-defect criterion.
Hence this ray is \(r(0,\phi)\), not merely P-maximal in test-curve
terminology.

On \(x\in(S_{i-1},S_i)\) the maximizing equation is
\(\rho=\log(x/(1-x))-t/i\). Relative to \(\omega\) the density is
\(e^{t/i}(1-x+e^{-t/i}x)^2\); on \((c,1)\) it is one. First derivatives
match at every interface, so there are no additional circle masses.
Integrating its logarithm against \(dx\) proves (5.6). The remainder is
between \(2\int_0^c\log(1-x)dx\geq-2\) and zero. Every finite-time
entropy is finite. This completes all assertions.

\textbf{Extension to finite positive real coefficients.} The same
statements hold for positive real \(a_i\) with \(c<1\), as actual smooth
relatively Kähler cohomological models. To check positivity, place the
coefficient vector in a small simplex with rational vertices, all still
satisfying strict (5.8). Each rational vertex gives a relatively ample
bundle as above. A positive convex combination of their Kähler
representatives, after one common large base twist, gives the required
real class. For compatibility no nonexistent real-exponent holomorphic
monomial is needed. On the local quadrant
\(\rho,\sigma=\log|\tau|^2\leq0\), \[
 h=\sum_i(a_i/i)\max(i\rho,\sigma).
\] Each actual ideal \((z^i,\tau)\) becomes invertible on this fan by
the same two-ray divisibility test. Its coefficient on \(F_k\) is
\(\min(i,k)\). Real weighted logarithms of their invertible local
factors give \(-d_k=\sum_i(a_i/i)\min(i,k)\). The comparison
\(|q-h|\leq\log2\) therefore proves real-divisor compatibility directly.
Off this quadrant ordinary smooth bundle frames suffice. The analytic
level and measure calculations depend only on real convexity and
intersections, so remain unchanged. This extension asserts cohomological
models, not rational line bundles when the coefficients are irrational.

\textbf{Theorem 5.2 (candidate188 reconstructed).} In Lemma 5.1 take
\(m=j\geq2\), \[
 a_i=j^{-3}\quad(i<j),\qquad
 a_j=1-j^{-1}-(j-1)j^{-3},\qquad c_j=1-j^{-1}.
\] All coefficients are positive rational numbers. Every datum has
supremum zero and norm \(r_j=a_j/j+j^{-3}\sum_{i<j}1/i\to0\) on all
semivaluations. The models are actually relatively ample, and their
energies and discrepancy entropies tend to zero. Uniform datum
convergence is strong convergence; directly the energy error is bounded
by \(\|g_j\|_\infty=r_j\). The actual entropy satisfies
\(-2\leq H(u_{j,t})-tr_j\leq0\), so every joint \(j,t\to\infty\) has
\(H(u_{j,t})/t\to0\).

Every fixed \(\lambda<0\) eventually satisfies \(\lambda\leq-r_j\),
giving \emph{exact} level zero. All cutoffs are exactly zero.
Nevertheless the endpoint is \(w_{1-1/j}\), with raw mass \(1/j\), polar
set \(\{0\}\), pole coefficient \(1-1/j\), and distance \((1-1/j)/2\)
from the target endpoint zero. The latter has mass one and empty polar
set. No moving-cutoff change repairs this example: the cutoffs already
coincide at zero. This refutes the implication from the listed recovery
properties to endpoint recovery. It does not refute existence: the
constant trivial model sequence for the target datum zero recovers its
endpoint exactly.

\subsection{6. A genuine endpoint-preserving construction for an
infinite
chain}\label{a-genuine-endpoint-preserving-construction-for-an-infinite-chain}

The next theorem is a new restricted construction. Its hypotheses are
explicit extra structure and are not imposed on the original target.

\textbf{Theorem 6.1 (finite tail with fixed mass and discrepancy
moment).} On \((\mathbf P^1,\omega)\) fix real \(a_i>0\), \(i\geq1\),
with \(c=\sum_{i\geq1}a_i\in(0,1)\), and put \(r=\sum a_i/i\). Define
\(S_i,g\) as in Lemma 5.1 using countably many intervals, with
\(S_i\uparrow c\). There exist actual smooth relatively Kähler
cohomological model data \(\phi_m\) converging uniformly and strongly to
\[ \phi(s\operatorname{ord}_0)=-\sum_{i\geq1}a_i\min(s,1/i)\] and zero
on other branches. Their energy converges and their discrepancy entropy
is \emph{exactly} \(r=H_\alpha(\phi)\) at every index. All their rays
start at zero, have cutoff zero, and have \emph{exactly the same
endpoint} \(w_c\) at every index. This includes its exact pole and raw
mass \(1-c\). Every fixed \(\lambda<0\) is eventually recovered exactly
as a P envelope, with exactly its mixed raw masses and pole coefficient.
Entropy divided by time tends to \(r\) in every joint limit. Their
entire level curves converge in \(d_{\mathcal S}\) uniformly in
\(\lambda\leq0\).

\emph{Proof of the construction.} Fix \(m\geq1\). Retain
\(a_1,\ldots,a_{m-1}\). Let
\[ T=\sum_{i\geq m}a_i,\qquad R=\sum_{i\geq m}a_i/i.\] Strict positivity
of infinitely many \(a_i\) gives \(0<R/T<1/m\). Choose a finite integer
\(N>m\) with \(1/N<R/T\). If \(N>m+1\), choose \(\varepsilon>0\)
sufficiently small, put \(\widetilde a_i=\varepsilon\) for \(m<i<N\),
and subtract these masses and moments: \[
 T'=T-\varepsilon(N-m-1),\qquad
 R'=R-\varepsilon\sum_{m<i<N}1/i.
\] Openness of the strict inequalities lets us arrange \(T'>0\) and
\(1/N<R'/T'<1/m\). If there are no intervening indices use
\(T'=T,R'=R\). Set \[
 \widetilde a_m=\frac{R'-T'/N}{1/m-1/N},\qquad
 \widetilde a_N=T'-\widetilde a_m.
 \tag{6.1}
\] Both are strictly positive. For \(i<m\) use \(\widetilde a_i=a_i\),
and omit indices \(i>N\). Every index \(1,\ldots,N\) now has a positive
coefficient, \[ \sum_{i=1}^N\widetilde a_i=c,\qquad
    \sum_{i=1}^N\widetilde a_i/i=r. \tag{6.2}\] Apply the
real-coefficient version of Lemma 5.1 on the actual smooth chain with
\(N\) blowups. Strict positivity on every component is thus retained;
zero-degree inserted components are not passed off as ample.

Let \(g_m\) be its speed function. On \([0,S_{m-1}]\), the original
prefix speeds and the remaining integral \(R\) agree, so
\[ g_m(x)=g(x)\quad(0\leq x\leq S_{m-1}). \tag{6.3}\] Both functions
vanish on \([c,1]\), and on the remaining tail their speeds are in
\([0,1/m]\). Hence \(\|g_m-g\|_\infty\leq T/m\to0\). Their dual datum
functions in (5.9) differ by at most this norm, on every finite
zero-branch valuation. At the supported zero semivaluation both equal
\(-r\) exactly; other supported branches have value zero. Uniform
convergence therefore holds on the original full carrier.

The limit is bounded psh NA, hence belongs to \(\mathcal E^{1,\rm NA}\).
Uniform energy continuity, or (5.5), gives
\(E_\alpha(\phi_m)\to-\int g=E_\alpha(\phi)\), so this convergence is
strong. Their model probability measures agree on the first \(m-1\)
atoms; each remaining part has mass \(T\), and the trivial atom has
weight \(1-c\) throughout. Their total variation difference from
\[ \nu=(1-c)\delta_{v_{\rm triv}}
              +\sum_{i\geq1}a_i\delta_{i^{-1}\operatorname{ord}_0}\] is
at most \(2T\to0\). Strong MA continuity identifies
\(\operatorname{MA}_\alpha(\phi)=\nu\). Its discrepancy integral is
\(r\); (6.2) gives exact entropy recovery, not just lower
semicontinuity.

Define \(q_t\) from the continuous convex limit \(g\) in (5.2). The
maximization bound gives
\(\|u_{m,t}-u_t\|_\infty\leq t\|g_m-g\|_\infty\), uniformly on bounded
time intervals. The infinite function \(g\) is strictly decreasing on
\([0,c)\), so (5.4) holds by the same convex minimization, including the
endpoint \(w_c\). These are P-model levels by the direct proof above and
form a bounded maximal test curve. The same analytic correspondence
makes \(u\) an actual finite-energy weak ray from zero. Alternatively,
joint maximal-ray continuity in R identifies the bounded-time \(d_1\)
limit of the already actual maximal rays with \(r(0,\phi)\). Thus the
datum here is not merely assigned to an unrelated test curve.

For \(-r<\lambda<0\), let \(b<c\) be the unique solution
\(g(b)=-\lambda\). For all sufficiently large \(m\), \(b<S_{m-1}\), so
(6.3) gives the exact same solution for \(g_m\). Equation (5.4) proves
eventual \emph{pointwise} envelope equality. Levels at or below \(-r\)
are zero at every index, and every endpoint is \(w_c\) at every index by
(6.2). Levels above zero are identically \(-\infty\) at every index.

For uniform control, on a fixed prefix \([0,S_{m-1}]\) the two level
maps are equal. For parameters whose inverse \(b\) lies in the tail,
both inverse values lie in \([S_{m-1},c]\), since the threshold value at
\(S_{m-1}\) agrees. Their difference is at most \(T=c-S_{m-1}\). Formula
(3.2) gives uniform \(d_{\mathcal S}\) error at most \(T/2\) on the
entire interval \(\lambda\leq0\), and uniform raw-mass error at most
\(T\). In dimension one these are all the mixed level masses (the other
mass is identically one).

The endpoint equality is exact at every index. Interior exactness is
eventual at each fixed parameter, not simultaneous equality of all
levels at any one finite index. Every nonzero fixed level has precisely
its original pole at zero once it is recovered; at the endpoint this
holds from the first index. No I-envelope replacement occurs.

Each finite model obeys (5.6) with \(r\) unchanged, so its joint entropy
limit is \(r\). The target ray also has finite-time entropy and slope
\(r\): the same density formula holds on the countably many \(x\)
intervals, whose endpoints have zero \(dx\) measure. Countable
integration gives \[ H(u_t)=tr+2\sum_{i\geq1}\int_{S_{i-1}}^{S_i}
                         \log(1-x+e^{-t/i}x)dx.\] The absolute remainder
is bounded by \(2\int_0^c-\log(1-x)dx\leq2\). There is no additional
curvature at the accumulating boundary \(c\): \(F+tg\) is strictly
convex, its dual derivative is continuous, and the one-sided maximizing
values tend to \(c\) on both sides. Its radial probability is still
\(dx\) over \([0,1]\). This completes the proof.

This is a constructive endpoint-preserving result for a specified
infinite-support datum. It uses extra one-dimensional chain structure,
initial zero, and a positive-mass endpoint \(c<1\). Its real
coefficients give real cohomological models, not generally rational
bundles. It proves no theorem for arbitrary log-discrepancy-model
inputs. Exact indexwise interior preservation of an arbitrary prescribed
continuum is also stronger than the conclusion just proved. Original
P19.2 remains open.

\subsection{7. Does a full-mass initial singularity leave all P ceilings
unchanged?}\label{does-a-full-mass-initial-singularity-leave-all-p-ceilings-unchanged}

There is a complete elementary statement for bounded starts. If
\(-C\leq w\leq0\), order in the initial point gives
\[ r(0,\phi)_t-C\leq r(w,\phi)_t\leq r(0,\phi)_t.\] Taking Legendre
infima gives the same two inequalities for every level. The definition
of \(P_\omega[\cdot]\) is invariant under bounded additive errors, so
all nontrivial absolute P ceilings agree, including endpoints.
Identically \(-\infty\) levels agree too. Exact polar sets and all raw
mixed masses agree because the potentials differ by a bounded amount.
The level functions themselves need not be equal, and this argument
proves no finite-time entropy continuity.

For an arbitrary full-mass \(\mathcal E^1\) start, blanket endpoint
invariance is false in a zero-mass regime. Here is a complete analytic
diagnostic, outside the finite-entropy-start branch of Theorem 4.2.

Use \(\zeta=1/z\) near infinity on \(\mathbf P^1\),
\(\rho=\log|\zeta|^2\), \(\delta=1/3\),
\(\rho_B=\log(\delta/(1-\delta))=-\log2\), and \(C=1-\delta-F(\delta)\).
Define a local total weight \[
 Q(\rho)=
 \begin{cases}
 -(-\rho)^{1/3}+C,&\rho\leq-1,\\
 \delta\rho-F(\delta),&-1\leq\rho\leq\rho_B,\\
 \log(1+e^\rho),&\rho\geq\rho_B,
 \end{cases}
 \qquad w=Q(\rho)-\log(1+e^\rho).
 \tag{7.1}
\] The values and first derivatives match. On the inner region
\(Q'=(1/3)(-\rho)^{-2/3}>0\) and \(Q''=(2/9)(-\rho)^{-5/3}>0\). These
derivatives join the affine piece of slope \(\delta\) and then the
convex Fubini--Study weight. Thus \(w\) is globally \(\omega\)-psh,
equals zero outside a disk about infinity, and is \(-\infty\) precisely
at infinity. It is nonpositive: inside \(\rho\leq-1\),
\(Q'(\rho)>e^\rho/(1+e^\rho)\), because \(L^{2/3}e^{-L}/(1+e^{-L})\) is
decreasing for \(L\geq1\) and its value at one is less than \(1/3\); the
difference increases to its nonpositive value at \(-1\). In the middle,
a convex tangent lies below the Fubini--Study function. Hence
\(\sup w=0\).

There is no polar atom, since \(Q'(-\infty)=0\). Its non-pluripolar
curvature has total mass \(Q'(+\infty)-Q'(-\infty)=1\). Moreover \[
 \int_{L\geq1}|w|\,dQ'\ \text{has tail bounded by a constant times}
             \int_1^\infty L^{1/3}L^{-5/3}dL<\infty.
\] The reference-volume integral of \(|w|\) is finite by its exponential
radial decay. On a curve these two integrals are exactly the finiteness
condition for \(E(w)=\tfrac12\int w(\omega+\omega_w)\). Thus
\(w\in\mathcal E^1(\omega)\) and has full mass despite its pole. Its
entropy is infinite: relative radial density on the inner region is
\(f=(2/9)L^{-5/3}/p'(-L)\) with \(p'(\rho)=e^\rho/(1+e^\rho)^2\). For
large \(L\), \(\log f\geq L/2\), and the positive entropy tail is at
least a constant times \(\int^\infty L^{-2/3}dL=+\infty\).

Now use the boundary nef model datum
\(\phi(s\operatorname{ord}_0)=-\min(s,1)\) and zero elsewhere. It
belongs to \(\mathcal E^{1,\rm NA}\) with energy \(-1/2\) and entropy
one, as in the supplied P19.1 calibration. Its zero-start actual ray is
\[ a_t=\log\frac{|z|^2+e^{-t}}{1+|z|^2},\qquad
 \chi_a(0)=w_1=\log|z|^2-\log(1+|z|^2).\] This nontrivial endpoint has
ordinary curvature \(\delta_0\) and raw non-pluripolar mass zero. Its P
property follows directly: any competitor below \(w_1+C\) has a
curvature atom of at least one at zero, hence its entire curvature is
\(\delta_0\), and it is \(w_1\) plus a constant. The constraint
\(h\leq0\) bounds that constant by zero. Therefore \(P_\omega[w_1]=w_1\)
even though the endpoint mass is zero. No positive-mass envelope theorem
is being applied here.

Let \(b=r(w,\phi)\). The constant ray \(w_t=w\) has actual usual datum
zero: its only local singularity grows as a cube root of a logarithm, so
on every modification its generic log-square divisorial coefficient is
zero. It is a finite-energy weak ray with energy slope zero; R's
energy-defect criterion makes it maximal for datum zero. Order in both
arguments now gives \(b_t\leq a_t\) and \(b_t\leq w\). For every point
where \(w\) is finite, the convex time function \(b_t\) is bounded
above, hence nonincreasing. At its pole the same inequalities force
\(b_t=-\infty\). Consequently \(h=\inf_t b_t\) is either an
\(\omega\)-psh function or identically \(-\infty\). If nontrivial,
\(h\leq w_1\) forces a curvature atom of at least one at zero. Total
curvature mass is one, so \(\omega_h=\delta_0\) and
\(h=w_1+\text{constant}\) on the compact curve. This is finite at
infinity, contradicting \(h\leq w\) there. Thus
\[ \chi_b(0)\equiv-\infty,\qquad
    d_1(a_t,b_t)=-E(w)<\infty\quad\text{for every }t. \tag{7.2}\] The
metric equality follows from order and (3.1), not from an envelope
theorem. Both zero-level raw masses are zero, but their P data differ.
This shows why mass alone and constant ray distance cannot identify an
arbitrary singular-start endpoint. This input has infinite initial
entropy and uses a boundary nef target, so it is not a counterexample to
the finite-entropy-start approximation theorem or the entire unchanged
P19.2 existence question. General unbounded finite-entropy-start P-level
invariance, including positive-mass and zero-mass endpoints, is not
proved by our entropy argument. We make no claim that every interior
ceiling changes in (7.2).

\subsection{8. Source applications, candidate checks, and unresolved
requirements}\label{source-applications-candidate-checks-and-unresolved-requirements}

Candidate187 is reconstructed by Theorem 4.2, using the permitted whole
P19.1 grant and the exact primary constant-invariant mixed estimate. Its
energy normalization must be applied before means are restored; using an
assumed singular supremum affinity or treating normalized \(d_1\) as
constant invariant would leave a gap. The reconstruction checks
finiteness at each index, uniform constants, the convex secant bound,
and the order of the fixed-test supremum. It proves rescaled joint
entropy, not fixed-time entropy convergence or P-endpoint preservation.

Candidate188 is reconstructed by Lemma 5.1 and Theorem 5.2. The proof
checks actual Cartier factors, every chart's divisibility, the node
\(x=0\), compatibility, positive degrees on every central component,
nonreduced fiber ampleness, the finite-dimensional Legendre
minimization, the P-envelope maximum principle, and the entropy
remainder. It is a valid diagnostic implication failure. The constant
trivial model still works for its target. No nonexistence inference for
the original target is justified.

The supplied complete136 paper and ledger were read as proof data and
retained locally. Its interior \(d_{\mathcal S}\) and mixed-mass results
are compatible with Theorem 5.2 and do not address its endpoint. They
are not used here as unconditional grants. We do not claim a new proof
of all their source chains. The supplied Cantor construction likewise is
not needed for the new constructive theorem or either reconstructed
candidate. If transported, its correct new potential is \(u_t=2v_{t/2}\)
with backgrounds unchanged: currents and probabilities agree,
\(E(u_t)=-a^2t/2\), ordinary datum remains zero, and entropy slope
remains infinite. Its energy defect \(a^2/2\) obstructs a particular
maximal-model mechanism, not the original finite-entropy prescribed
P-endpoint existence problem. No claim in this paper is inferred from
its historical review verdict.

The resulting new constructive signature is the explicit real
infinite-chain theorem, with exact endpoint, strong energy convergence,
exact discrepancy entropy and eventual exact fixed interior levels. The
following gaps remain for unchanged P19.2:

\begin{itemize}
\tightlist
\item
  A simultaneous construction, or genuine nonexistence theorem, for
  arbitrary residual log-discrepancy-model and prescribed P data,
  without the special infinite-chain structure, support assumptions, or
  positive endpoint hypothesis of Theorem 6.1.
\item
  A justified general endpoint topology and endpoint-envelope
  convergence, including zero-mass levels; all mixed masses do not
  determine a nontrivial zero-mass endpoint.
\item
  Exact indexwise prescribed singularities and poles where those are
  required. Eventual equality at each fixed interior level is a distinct
  quantifier from equality of an entire curve at every index.
\item
  Relative P-level assertions for arbitrary unbounded finite-entropy
  initial points, separate from their proved joint entropy limit.
\item
  Any fixed-time entropy convergence or genuinely hybrid measure
  statement required beyond the radial entropy conclusions proved here.
\end{itemize}

The original dossier does not become a formal universal endpoint theorem
merely by interpreting these gaps as definitions. We report this exact
limited progress and leave the unchanged residual atom PARTIAL. No
degenerating-family compactness (P19.3), automorphism or singular
stability extension (P19.4), or new completion of P19.1 is claimed.

\subsection{9. Reproducibility and evidence
boundaries}\label{reproducibility-and-evidence-boundaries}

The frozen \texttt{input.json} and \texttt{input-original.json} are
byte-identical, with SHA256:

\begin{Shaded}
\begin{Highlighting}[]
\NormalTok{3ad4e22fbef13e9e8f7ec9287b874397bf458fd880bc48d6fefef37d8cc7d0a6}
\end{Highlighting}
\end{Shaded}

\texttt{assignment-original.txt} is its exact decoded assignment and
\texttt{task-original.json} its decoded task object. The extraction
script preserves the complete supplied mathematical units in
\texttt{supplied/}. No other campaign, repository, database, transcript,
worker, or model was accessed. The existing execution model was not
changed.

Fresh public source tool responses and their line indices are saved in
\texttt{sources/}; exact statements, proof units actually read,
applicability, and limitations are recorded in
\texttt{source-application-ledger.json}. The only shell PDF attempt
failed because the inherited proxy was unreachable; no proxy, CA,
authentication, account or sandbox was changed. Public web-tool text
supplied the primary material. We claim no access to missing PDF image
pages and no full recursive reproof of standard source foundations. No
supervised heavy-compute facility was exposed; none was needed or
launched.

\texttt{check\_models.py} exhaustively checks rational chain identities
and monomial divisibility for its stated finite domain, and separately
records floating-point convex and entropy observations. These do not
prove the analytic theorems or replace the arguments above. It also
implements a rational instance of the two-moment tail construction. Run
\texttt{python\ check\_models.py}, \texttt{python\ build\_paper.py}, and
\texttt{python\ validate\_deliverables.py}. The build generates
standalone \texttt{P19.2-paper.tex}, runs existing pdfLaTeX twice, and
retains both command outputs, exit statuses, and complete logs.
Execution success is not a mathematical proof certificate.
\texttt{status.json} has exactly the five requested keys and records
PARTIAL.

\subsection{References}\label{references}

\begin{itemize}
\tightlist
\item
  {[}MP{]} P. Mesquita--Piccione, \emph{A non-Archimedean theory of
  complex spaces and the cscK problem},
  \href{https://arxiv.org/pdf/2409.06221v2}{arXiv:2409.06221v2}, Lemma
  5.3.2 and Theorem 5.3.4.
\item
  {[}MPWN{]} P. Mesquita--Piccione and D. Witt Nyström, \emph{A
  transcendental non-Archimedean Calabi--Yau theorem with applications
  to the cscK problem},
  \href{https://arxiv.org/pdf/2509.09442v1}{arXiv:2509.09442v1}, Lemma
  5.0.5 and the finite-energy strong topology.
\item
  {[}DZ25{]} T. Darvas and K. Zhang, \emph{A YTD correspondence for
  constant scalar curvature metrics},
  \href{https://arxiv.org/pdf/2509.15173v2}{arXiv:2509.15173v2}, formula
  (11), Section 9, Theorem 10.1. Used for exact scope comparison.
\item
  {[}BDL{]} R. Berman, T. Darvas and C. H. Lu, \emph{Convexity of the
  extended K-energy and the large time behavior of the weak Calabi
  flow}, Geometry \& Topology 21 (2017), 2945--2988,
  \href{https://msp.org/gt/2017/21-5/gt-v21-n5-p08-s.pdf}{published
  primary text}, Lemma 3.1 and Theorem 4.7.
\item
  {[}DZ22{]} T. Darvas and K. Zhang,
  \href{https://arxiv.org/pdf/2208.08324v4}{arXiv:2208.08324v4}, Theorem
  3.7(ii) and the analytic Legendre/Perron proof.
\item
  {[}DXZ{]} T. Darvas, M. Xia and K. Zhang, \emph{A transcendental
  approach to non-Archimedean metrics of pseudoeffective classes},
  Commentarii Mathematici Helvetici 100 (2025), 269--322,
  \href{https://ems.press/content/serial-article-files/50355}{published
  primary text} and \href{https://arxiv.org/pdf/2302.02541v3}{arXiv v3},
  Section 2.2, Theorem 2.6. Used to check relative signature; not to
  deduce endpoint invariance.
\item
  {[}DDNL{]} T. Darvas, E. Di Nezza and C. H. Lu, \emph{The metric
  geometry of singularity types},
  \href{https://arxiv.org/pdf/1909.00839}{arXiv:1909.00839}, Lemma 3.4
  and its preceding chordal-energy formula.
\item
  {[}DDNL-products{]} T. Darvas, E. Di Nezza and C. H. Lu,
  \emph{Monotonicity of non-pluripolar products and complex
  Monge--Ampère equations with prescribed singularity},
  \href{https://arxiv.org/pdf/1705.05796v3}{arXiv:1705.05796v3}, Theorem
  2.4.
\item
  {[}Stacks{]} The Stacks Project,
  \href{https://stacks.math.columbia.edu/tag/0B5Y}{0B5Y},
  \href{https://stacks.math.columbia.edu/tag/09MS}{09MS}, and
  \href{https://stacks.math.columbia.edu/tag/0D2S}{0D2S}.
\item
  {[}R{]} Complete supplied P19.1 mathematical paper, \emph{Entropy
  slopes of finite-energy Kähler rays: calibration, recovery, and
  arbitrary starting points}, 6 October 2026. Used only as the permitted
  theorem grant, with its complete retained mathematical signature.
\end{itemize}

\end{document}
